\documentclass[10pt,a4paper]{amsart}
\usepackage{amsmath,amssymb,amsthm,graphicx}
% Retained preamble for the Stacks Project chapter snapshot in this directory.
% The chapter's own preamble.tex is a copy of the shared Stacks Project preamble; the block
% below is sliced MECHANICALLY, line for line, from it. Nothing was added, removed or
% redefined; the wrapper class, the xr-hyper/hyperref wiring and the externaldocument list
% are omitted because the assembler supplies the document class.
%
% lines taken from preamble.tex: 3, 5, 9, 11, 12, 156, 157, 158, ... (51 lines total)
%
% package coverage restored from the chapter's own preamble (2026-10-13): xy, xy-2cell

\usepackage[all]{xy}
\xyoption{2cell}
\UseAllTwocells

\newenvironment{reference}{\comment}{\endcomment}
\newenvironment{slogan}{\comment}{\endcomment}
\newenvironment{history}{\comment}{\endcomment}
\theoremstyle{plain}
\newtheorem{theorem}[subsection]{Theorem}
\newtheorem{proposition}[subsection]{Proposition}
\newtheorem{lemma}[subsection]{Lemma}
\theoremstyle{definition}
\newtheorem{definition}[subsection]{Definition}
\newtheorem{example}[subsection]{Example}
\newtheorem{exercise}[subsection]{Exercise}
\newtheorem{situation}[subsection]{Situation}
\theoremstyle{remark}
\newtheorem{remark}[subsection]{Remark}
\newtheorem{remarks}[subsection]{Remarks}
\numberwithin{equation}{subsection}
\def\lim{\mathop{\mathrm{lim}}\nolimits}
\def\colim{\mathop{\mathrm{colim}}\nolimits}
\def\Spec{\mathop{\mathrm{Spec}}}
\def\Hom{\mathop{\mathrm{Hom}}\nolimits}
\def\Ext{\mathop{\mathrm{Ext}}\nolimits}
\def\SheafHom{\mathop{\mathcal{H}\!\mathit{om}}\nolimits}
\def\SheafExt{\mathop{\mathcal{E}\!\mathit{xt}}\nolimits}
\def\Sch{\mathit{Sch}}
\def\Mor{\mathop{\mathrm{Mor}}\nolimits}
\def\Ob{\mathop{\mathrm{Ob}}\nolimits}
\def\Sh{\mathop{\mathit{Sh}}\nolimits}
\def\NL{\mathop{N\!L}\nolimits}
\def\CH{\mathop{\mathrm{CH}}\nolimits}
\def\proetale{{pro\text{-}\acute{e}tale}}
\def\etale{{\acute{e}tale}}
\def\QCoh{\mathit{QCoh}}
\def\Ker{\mathop{\mathrm{Ker}}}
\def\Im{\mathop{\mathrm{Im}}}
\def\Coker{\mathop{\mathrm{Coker}}}
\def\Coim{\mathop{\mathrm{Coim}}}
\DeclareMathSymbol{\boxtimes}{\mathbin}{AMSa}{"02}
\def\QCohstack{\mathcal{QC}\!\mathit{oh}}
\def\Cohstack{\mathcal{C}\!\mathit{oh}}
\def\Spacesstack{\mathcal{S}\!\mathit{paces}}
\def\Quotfunctor{\mathrm{Quot}}
\def\Hilbfunctor{\mathrm{Hilb}}
\def\Curvesstack{\mathcal{C}\!\mathit{urves}}
\def\Polarizedstack{\mathcal{P}\!\mathit{olarized}}
\def\Complexesstack{\mathcal{C}\!\mathit{omplexes}}
\def\Pic{\mathop{\mathrm{Pic}}\nolimits}
\def\Picardstack{\mathcal{P}\!\mathit{ic}}
\def\Picardfunctor{\mathrm{Pic}}
\def\Deformationcategory{\mathcal{D}\!\mathit{ef}}
\usepackage{cleveref}
\crefname{originalproposition}{original proposition}{original propositions}
\Crefname{originalproposition}{Original proposition}{Original propositions}
\makeatletter
\newlabel{categories-lemma-yoneda}{{001P}{source}{}{sourcepointer-1}{}}
\newlabel{categories-lemma-yoneda@cref}{{[originalproposition][0][]001P}{[1][0][]source}{}{}{}}
\makeatother
\crefname{originalproposition}{original proposition}{original propositions}
\Crefname{originalproposition}{Original proposition}{Original propositions}
\makeatletter
\newlabel{injectives-lemma-grothendieck-brown}{{07D7}{source}{}{sourcepointer-2}{}}
\newlabel{injectives-lemma-grothendieck-brown@cref}{{[originalproposition][0][]07D7}{[1][0][]source}{}{}{}}
\makeatother
\crefname{originalproposition}{original proposition}{original propositions}
\Crefname{originalproposition}{Original proposition}{Original propositions}
\makeatletter
\newlabel{injectives-lemma-grothendieck-products}{{07D8}{source}{}{sourcepointer-3}{}}
\newlabel{injectives-lemma-grothendieck-products@cref}{{[originalproposition][0][]07D8}{[1][0][]source}{}{}{}}
\makeatother
\crefname{originalproposition}{original proposition}{original propositions}
\Crefname{originalproposition}{Original proposition}{Original propositions}
\makeatletter
\newlabel{injectives-section-grothendieck-conditions}{{079A}{source}{}{sourcepointer-4}{}}
\newlabel{injectives-section-grothendieck-conditions@cref}{{[originalproposition][0][]079A}{[1][0][]source}{}{}{}}
\makeatother
\crefname{originalproposition}{original proposition}{original propositions}
\Crefname{originalproposition}{Original proposition}{Original propositions}
\makeatletter
\newlabel{injectives-theorem-injective-embedding-grothendieck}{{079H}{source}{}{sourcepointer-5}{}}
\newlabel{injectives-theorem-injective-embedding-grothendieck@cref}{{[originalproposition][0][]079H}{[1][0][]source}{}{}{}}
\makeatother
\crefname{originalproposition}{original proposition}{original propositions}
\Crefname{originalproposition}{Original proposition}{Original propositions}
\makeatletter
\newlabel{lemma-quasi-coherent-algebraic-stack}{{lemma-quasi-coherent-algebraic-stack}{source}{}{sourcepointer-6}{}}
\newlabel{lemma-quasi-coherent-algebraic-stack@cref}{{[originalproposition][0][]lemma-quasi-coherent-algebraic-stack}{[1][0][]source}{}{}{}}
\makeatother
\crefname{originalproposition}{original proposition}{original propositions}
\Crefname{originalproposition}{Original proposition}{Original propositions}
\makeatletter
\newlabel{properties-proposition-coherator}{{077P}{source}{}{sourcepointer-7}{}}
\newlabel{properties-proposition-coherator@cref}{{[originalproposition][0][]077P}{[1][0][]source}{}{}{}}
\makeatother
\crefname{originalproposition}{original proposition}{original propositions}
\Crefname{originalproposition}{Original proposition}{Original propositions}
\makeatletter
\newlabel{sites-modules-lemma-limits-colimits}{{03DB}{source}{}{sourcepointer-8}{}}
\newlabel{sites-modules-lemma-limits-colimits@cref}{{[originalproposition][0][]03DB}{[1][0][]source}{}{}{}}
\makeatother
\crefname{originalproposition}{original proposition}{original propositions}
\Crefname{originalproposition}{Original proposition}{Original propositions}
\makeatletter
\newlabel{spaces-groupoids-lemma-set-generators}{{0780}{source}{}{sourcepointer-9}{}}
\newlabel{spaces-groupoids-lemma-set-generators@cref}{{[originalproposition][0][]0780}{[1][0][]source}{}{}{}}
\makeatother
\crefname{originalproposition}{original proposition}{original propositions}
\Crefname{originalproposition}{Original proposition}{Original propositions}
\makeatletter
\newlabel{spaces-properties-proposition-coherator}{{077V}{source}{}{sourcepointer-10}{}}
\newlabel{spaces-properties-proposition-coherator@cref}{{[originalproposition][0][]077V}{[1][0][]source}{}{}{}}
\makeatother
\title{Stacks, Sheaves on Algebraic Stacks: proposition-coherator}
\author{The Stacks Project Authors}
\date{stacks-project:stacks-sheaves.tex:local-snapshot-20261012; GNU FDL 1.2 or later}
\begin{document}
\maketitle
\noindent\textit{Editorial note.} Statement span is lines 2656-2676 of sources/stacks-algebraic-stacks-sheaves/pinned-original-source.tex; proof span is lines 2678-2748, the complete author proof (69 body lines).\par\medskip
\noindent\textit{Editorial note.} Navigation stubs: 10 cited-but-not-included labels are supplied as non-load-bearing pointers with their pinned tags where the snapshot tag index has them.\par\medskip
\noindent\textit{Editorial note.} extra\_packages is empty on purpose: the retained preamble already loads xy with options and a second declaration clashes in the assembler wrapper.\par\medskip
\noindent\textit{Editorial note.} Text is retained verbatim, including the author's own typos and placeholders.\par\medskip
\noindent\textit{Editorial note.} Queue index 6 of queue\_algebraic-stacks-sheaves.json (pb=2918, stmt 2656-2676, proof 2678-2748).\par\medskip
\noindent\textit{Editorial note.} Cross-chapter same-label warning (NOT an occupancy conflict, Chief ruling 2026-10-13 evening): label 'proposition-coherator' also exists in stacks-spaces-properties — a different proposition in a different chapter file.\par\medskip
\noindent\textit{Allowed standard prerequisites (editorial).}\begin{itemize}
\item Explicit prerequisite boundary. The cited inputs are prerequisites; the argument in the retained span is not replaced by any cited result.
\end{itemize}\medskip
\ifdefined\hypertarget\hypertarget{sourcepointer-1}{}\fi
\noindent\textit{Source reference.} Original proposition 001P: sources/stacks-algebraic-stacks-sheaves/pinned-original-source.tex, source lines 2748--2748. Cited by the retained statement or proof but not part of the collected spans; supplied only as a navigation pointer. The argument does not depend on its content.\par\medskip
\ifdefined\hypertarget\hypertarget{sourcepointer-2}{}\fi
\noindent\textit{Source reference.} Original proposition 07D7: sources/stacks-algebraic-stacks-sheaves/pinned-original-source.tex, source lines 2748--2748. Cited by the retained statement or proof but not part of the collected spans; supplied only as a navigation pointer. The argument does not depend on its content.\par\medskip
\ifdefined\hypertarget\hypertarget{sourcepointer-3}{}\fi
\noindent\textit{Source reference.} Original proposition 07D8: sources/stacks-algebraic-stacks-sheaves/pinned-original-source.tex, source lines 2748--2748. Cited by the retained statement or proof but not part of the collected spans; supplied only as a navigation pointer. The argument does not depend on its content.\par\medskip
\ifdefined\hypertarget\hypertarget{sourcepointer-4}{}\fi
\noindent\textit{Source reference.} Original proposition 079A: sources/stacks-algebraic-stacks-sheaves/pinned-original-source.tex, source lines 2748--2748. Cited by the retained statement or proof but not part of the collected spans; supplied only as a navigation pointer. The argument does not depend on its content.\par\medskip
\ifdefined\hypertarget\hypertarget{sourcepointer-5}{}\fi
\noindent\textit{Source reference.} Original proposition 079H: sources/stacks-algebraic-stacks-sheaves/pinned-original-source.tex, source lines 2748--2748. Cited by the retained statement or proof but not part of the collected spans; supplied only as a navigation pointer. The argument does not depend on its content.\par\medskip
\ifdefined\hypertarget\hypertarget{sourcepointer-6}{}\fi
\noindent\textit{Source reference.} Original proposition lemma-quasi-coherent-algebraic-stack: sources/stacks-algebraic-stacks-sheaves/pinned-original-source.tex, source lines 2601--2631. Cited by the retained statement or proof but not part of the collected spans; supplied only as a navigation pointer. The argument does not depend on its content.\par\medskip
\ifdefined\hypertarget\hypertarget{sourcepointer-7}{}\fi
\noindent\textit{Source reference.} Original proposition 077P: sources/stacks-algebraic-stacks-sheaves/pinned-original-source.tex, source lines 2748--2748. Cited by the retained statement or proof but not part of the collected spans; supplied only as a navigation pointer. The argument does not depend on its content.\par\medskip
\ifdefined\hypertarget\hypertarget{sourcepointer-8}{}\fi
\noindent\textit{Source reference.} Original proposition 03DB: sources/stacks-algebraic-stacks-sheaves/pinned-original-source.tex, source lines 2748--2748. Cited by the retained statement or proof but not part of the collected spans; supplied only as a navigation pointer. The argument does not depend on its content.\par\medskip
\ifdefined\hypertarget\hypertarget{sourcepointer-9}{}\fi
\noindent\textit{Source reference.} Original proposition 0780: sources/stacks-algebraic-stacks-sheaves/pinned-original-source.tex, source lines 2748--2748. Cited by the retained statement or proof but not part of the collected spans; supplied only as a navigation pointer. The argument does not depend on its content.\par\medskip
\ifdefined\hypertarget\hypertarget{sourcepointer-10}{}\fi
\noindent\textit{Source reference.} Original proposition 077V: sources/stacks-algebraic-stacks-sheaves/pinned-original-source.tex, source lines 2748--2748. Cited by the retained statement or proof but not part of the collected spans; supplied only as a navigation pointer. The argument does not depend on its content.\par\medskip
\begin{proposition}
\label{proposition-coherator}
Let $\mathcal{X}$ be an algebraic stack over $S$.
\begin{enumerate}
\item The category $\QCoh(\mathcal{O}_\mathcal{X})$ is a Grothendieck
abelian category. Consequently, $\QCoh(\mathcal{O}_\mathcal{X})$
has enough injectives and all limits.
\item The inclusion functor
$\QCoh(\mathcal{O}_\mathcal{X}) \to
\textit{Mod}(\mathcal{O}_\mathcal{X})$ has a right adjoint\footnote{This
functor is sometimes called the {\it coherator}.}
$$
Q :
\textit{Mod}(\mathcal{O}_\mathcal{X})
\longrightarrow
\QCoh(\mathcal{O}_\mathcal{X})
$$
such that for every quasi-coherent sheaf $\mathcal{F}$ the adjunction mapping
$Q(\mathcal{F}) \to \mathcal{F}$ is an isomorphism.
\end{enumerate}
\end{proposition}

\begin{proof}
This proof is a repeat of the proof in the case of schemes, see
Properties, Proposition \ref{properties-proposition-coherator}
and the case of algebraic spaces, see
Properties of Spaces, Proposition
\ref{spaces-properties-proposition-coherator}.
We advise the reader to read either of those proofs first.

\medskip\noindent
Part (1) means $\QCoh(\mathcal{O}_\mathcal{X})$ (a) has all colimits,
(b) filtered colimits are exact, and (c) has a generator, see
Injectives, Section \ref{injectives-section-grothendieck-conditions}.
By Lemma \ref{lemma-quasi-coherent-algebraic-stack}
colimits in $\QCoh(\mathcal{O}_X)$ exist and agree
with colimits in $\textit{Mod}(\mathcal{O}_X)$. By
Modules on Sites, Lemma \ref{sites-modules-lemma-limits-colimits}
filtered colimits are exact. Hence (a) and (b) hold.

\medskip\noindent
Choose a presentation $\mathcal{X} = [U/R]$ so that $(U, R, s, t, c)$
is a smooth groupoid in algebraic spaces and in particular $s$ and $t$
are flat morphisms of algebraic spaces. By
Lemma \ref{lemma-quasi-coherent-algebraic-stack}
above we have
$\QCoh(\mathcal{O}_\mathcal{X}) = \QCoh(U, R, s, t, c)$.
By Groupoids in Spaces, Lemma \ref{spaces-groupoids-lemma-set-generators}
there exists a set $T$ and a family $(\mathcal{F}_t)_{t \in T}$ of
quasi-coherent sheaves on $\mathcal{X}$ such that every quasi-coherent
sheaf on $\mathcal{X}$ is the directed colimit of its subsheaves
which are isomorphic to one of the $\mathcal{F}_t$.
Thus $\bigoplus_t \mathcal{F}_t$ is
a generator of $\QCoh(\mathcal{O}_X)$ and we conclude that (c) holds.
The assertions on limits and injectives hold in any
Grothendieck abelian category, see
Injectives, Theorem
\ref{injectives-theorem-injective-embedding-grothendieck} and
Lemma \ref{injectives-lemma-grothendieck-products}.

\medskip\noindent
Proof of (2). To construct $Q$ we use the following general procedure.
Given an object $\mathcal{F}$ of $\textit{Mod}(\mathcal{O}_\mathcal{X})$
we consider the functor
$$
\QCoh(\mathcal{O}_\mathcal{X})^{opp}
\longrightarrow
\textit{Sets},
\quad
\mathcal{G}
\longmapsto
\Hom_\mathcal{X}(\mathcal{G}, \mathcal{F})
$$
This functor transforms colimits into limits,
hence is representable, see
Injectives, Lemma \ref{injectives-lemma-grothendieck-brown}.
Thus there exists a quasi-coherent sheaf $Q(\mathcal{F})$
and a functorial isomorphism
$\Hom_\mathcal{X}(\mathcal{G}, \mathcal{F}) =
\Hom_\mathcal{X}(\mathcal{G}, Q(\mathcal{F}))$
for $\mathcal{G}$ in $\QCoh(\mathcal{O}_\mathcal{X})$.
By the Yoneda lemma
(Categories, Lemma \ref{categories-lemma-yoneda})
the construction $\mathcal{F} \leadsto Q(\mathcal{F})$ is
functorial in $\mathcal{F}$. By construction $Q$ is a right
adjoint to the inclusion functor.
The fact that $Q(\mathcal{F}) \to \mathcal{F}$ is an isomorphism
when $\mathcal{F}$ is quasi-coherent is a formal consequence of the fact
that the inclusion functor
$\QCoh(\mathcal{O}_\mathcal{X}) \to
\textit{Mod}(\mathcal{O}_\mathcal{X})$
is fully faithful.
\end{proof}

\end{document}
